\chapter{Memory Loss}

\section*{Definition}
Impairment in ability to acquire new information or recall previously learned information.

\section*{Pathophysiology}
Distinguishes between normal age-related forgetting and pathological impairment such as amnestic syndromes dementia or transient global amnesia; episodic memory most commonly affected

Alzheimer's disease (p.170) progressive neurodegenerative disorder characterized by gradual onset memory impairment affecting recent recall first then extends to remote memory accompanied by cognitive decline in language orientation executive function visuospatial skills behavioral changes mood disturbances agitation apathy depression psychosis sundowning syndrome functional independence progressively impaired requiring increasing care support eventual complete dependence end-stage bedbound incontinence nonverbal pathological hallmarks amyloid beta plaques neurofibrillary tangles tau protein neuronal loss cerebral atrophy hippocampal entorhinal cortex medial temporal lobe involvement early stages frontal parietal association areas later widespread cortical atrophy diagnosis clinical evaluation supported neuropsychological testing biomarkers CSF amyloid tau PET imaging APOE genetic risk factor no cure management cholinesterase inhibitors memantine supportive care safety modifications environmental adaptations caregiver support advance directives planning

Vascular dementia stepwise progression correlated with cerebrovascular events strategic infarcts multi-infarct state white matter ischemic small vessel disease subcortical vascular cognitive impairment depression pseudo-differential features from Alzheimer typically earlier gait disturbance urinary incontinence focal neurological signs pseudobulbar affect

Depression (pseudodementia) subjective complaints of memory difficulties often exaggerated apparent slowing of response times on examination concentration effort poor motivation depressive symptoms prominent true objective memory testing reveals intact recall with prompting improves with encouragement treatment responsive reversible unlike organic dementia

Thyroid dysfunction hypothyroidism can cause cognitive slowing memory complaints fatigue apathy weight gain cold dry skin bradycardia deep tendon reflexes delayed relaxation phase reversible with thyroid hormone replacement

Medication side effects anticholinergics benzodiazepines opioids antihistamines anticonvulsants alcohol sedatives hypnotics may impair attention encoding retrieval processes polypharmacy elderly particularly vulnerable deprescribing review medication reconciliation essential component assessment

\section*{Etiology}
\begin{itemize}
  \item Alzheimer disease progressive neurodegeneration memory decline language dysfunction executive impairment behavioral changes eventually complete dependence
  \item Vascular dementia stepwise cognitive decline associated with cerebrovascular events gait disturbance urinary incontinence focal signs
  \item Depression pseudodementia subjective memory complaints concentration effort poor motivation improves with treatment realistic recall with prompting
  \item Thyroid dysfunction hypothyroidism cognitive slowing fatigue weight gain cold intolerance dry skin bradycardia reversible with therapy
  \item Medication side effects anticholinergics benzodiazepines opioids sedatives hypnotics alcohol especially problematic polypharmacy elderly patients
\end{itemize}

\section*{Indications for Evaluation}
Seek care if:
\begin{itemize}
  \item Severe or worsening symptoms
  \item Symptoms lasting more than progressive or sudden onset warrants neurological evaluation including imaging and neuropsychological testing
  \item Accompanied by other concerning signs personality change gait disturbance incontinence focal deficits seizures
\end{itemize}

\section*{Related Symptoms}
\begin{itemize}
  \item Getting lost easily Repetitive questions Difficulty learning new names Misplacing items frequently Confusion about dates/times Executive dysfunction Language problems Visuospatial difficulties Mood changes Behavioral changes Sleep disturbances Appetite changes
\end{itemize}
\textbf{Note:} Always consider serious underlying conditions when evaluating persistent memory loss.

\section*{Self-Care / Home Management}
\begin{itemize}
  \item Ensure a safe environment to prevent injury during episodes (remove hazards, avoid driving or operating machinery).
  \item Keep a symptom diary documenting triggers, duration, frequency, and associated features.
  \item Practice stress reduction techniques (deep breathing, meditation, adequate sleep) as stress often exacerbates neurological symptoms.
  \item Avoid alcohol, caffeine, and recreational drugs which can worsen many neurological conditions.
  \item Seek immediate evaluation for sudden-onset, severe, or progressive neurological symptoms.
\end{itemize}

\section*{Diagnostic Approach}
\begin{itemize}
  \item Detailed neurological history includes onset pattern, progression, triggers, associated symptoms, and functional impact.
  \item Comprehensive neurological examination assesses mental status, cranial nerves, motor/sensory function, reflexes, coordination, and gait.
  \item Neuroimaging (CT, MRI) is often indicated to evaluate for structural, vascular, or demyelinating causes.
  \item Specialized testing (EEG for seizures, nerve conduction studies for neuropathy, lumbar puncture for meningitis) is guided by clinical suspicion.
\end{itemize}

\section*{Treatment / Management Overview}
\begin{itemize}
  \item Treatment is directed at the underlying cause identified through diagnostic evaluation.
  \item Symptomatic pharmacotherapy may include analgesics, anticonvulsants, antidepressants, or disease-modifying agents as appropriate.
  \item Physical, occupational, and speech therapy play important roles in functional recovery and rehabilitation.
  \item Referral to neurology is indicated for unexplained, progressive, or treatment-refractory neurological symptoms.
\end{itemize}

\clearpage
